\section{Validity}

评估是否有效，是比”分数是多少”更底层的问题。

\subsection{Train-test overlap}

机器学习最基本的原则是不要在测试集上训练。可是在互联网语料时代，训练数据来自整个网络，是否污染很难彻底判断。

课程提到两条思路：

\begin{itemize}
    \item \textbf{从模型侧推断污染}：利用数据点之间的可交换性等性质去估计重叠。
    \item \textbf{从流程侧要求披露}：让模型提供方报告 train-test overlap 或类似统计。
\end{itemize}

\begin{warningbox}{Benchmark Contamination：数据泄露的系统性威胁}
Benchmark contamination（基准污染）是当前大模型评估中最严重的系统性问题之一。由于预训练数据来自整个互联网，而许多 benchmark 的题目和答案也公开在网上，模型很可能在训练阶段就``见过''测试题。

具体的污染路径包括：
\begin{itemize}
    \item \textbf{直接泄露}：benchmark 的题目和答案被收录进训练语料（如 Wikipedia 引用、学术论坛讨论、GitHub issue）。
    \item \textbf{间接泄露}：训练数据中包含关于 benchmark 的讨论、解题教程或排行榜分析。
    \item \textbf{合成数据放大}：用其他模型生成训练数据时，如果那些模型的训练也包含了 benchmark 数据，污染就通过合成数据间接传递。
    \item \textbf{评估集 rephrasing}：有些团队把评估题稍作改写后放入训练集，形式上不算直接泄露，但模型仍然能从中获益。
\end{itemize}

检测 contamination 的方法包括：检查 n-gram overlap、用 membership inference 判断模型是否``记住了''某个数据点、以及比较模型在公开题与未公开变体上的表现差异。但没有一种方法能做到完全可靠。
\end{warningbox}

\subsection{Goodhart's Law 与过拟合 Benchmark}

\begin{warningbox}{Goodhart's Law：当 benchmark 变成目标，它就不再是好的度量}
``When a measure becomes a target, it ceases to be a good measure.'' 这条经济学定律在 LLM 评估领域表现得尤为突出。

具体表现包括：
\begin{itemize}
    \item \textbf{训练数据选择偏向}：如果知道模型会被 MMLU 评估，就有动机在训练数据中增加教科书和考试题的比例。
    \item \textbf{prompt engineering 过拟合}：为特定 benchmark 的 prompt 格式进行专门优化，使模型在该格式下表现远好于其他格式。
    \item \textbf{reward hacking}：在 RLHF 阶段，如果 reward model 与某些评估指标高度相关，模型会学会``讨好''评估指标而非真正提升能力。
    \item \textbf{系统级优化}：在 agent benchmark 中，团队可能为特定 benchmark 编写专门的 scaffolding 代码，使得分数反映的是工程投入而非模型能力。
\end{itemize}

Goodhart's Law 的深层含义是：评估永远只是真实目标的近似。一旦近似本身成为优化目标，它和真实目标的差距就会被放大。这也是为什么评估社区需要不断发布新的 benchmark：不仅因为旧的被饱和了，也因为旧的被优化废了。
\end{warningbox}

\subsection{Dataset quality}

不是所有 benchmark 都是一开始就足够干净。比如 SWE-bench Verified 就是对原始 SWE-bench 做进一步修正得到的更可靠版本。更一般地，很多 benchmark 都需要通过人工或半自动方式不断做 ``platinum'' 级别的清洗。

数据集质量问题的常见类型：

\begin{table}[H]
\centering
\begin{tabular}{p{3.0cm}p{4.5cm}p{5.5cm}}
\toprule
\textbf{质量问题} & \textbf{具体表现} & \textbf{对评估的影响} \\
\midrule
标签错误 & 参考答案本身就是错的 & 模型答对反而被判错，拉低得分 \\
题目歧义 & 多个选项都可以被合理论证为正确 & 不同评估 prompt 给出不同结果 \\
难度分布偏斜 & 大部分题太简单或太难 & 无法区分中等水平的模型 \\
领域覆盖不均 & 某些子领域题目远多于其他 & 宏观分数被高频领域主导 \\
时效性过期 & 题目基于过时的事实 & 模型学了更新的知识反而答``错'' \\
\bottomrule
\end{tabular}
\caption{常见的 benchmark 数据集质量问题及其影响}
\end{table}

\begin{knowledgebox}{评估有效性与数据质量}
如果 benchmark 自身有错误、歧义或污染，那么模型排名的细微差异可能没有你想的那么可信。很多时候，改进 benchmark 和改进模型同样重要。
\end{knowledgebox}

\subsection{统计显著性：分差到底够不够大}

即使 benchmark 数据质量没有问题，一个常被忽略的问题是：两个模型之间的分数差异是否具有统计显著性？

如果一个 benchmark 有 $n$ 道题，模型准确率为 $p$，那么准确率的标准误差近似为：
\[
\mathrm{SE} = \sqrt{\frac{p(1-p)}{n}}
\]

以 MMLU 为例，如果模型准确率约为 85\%，题目数约为 14000，则：
\[
\mathrm{SE} \approx \sqrt{\frac{0.85 \times 0.15}{14000}} \approx 0.003
\]

这意味着 95\% 置信区间大约是 $\pm 0.6\%$。因此，两个模型如果分差不到 1\%，在统计上可能无法区分。但在 MMLU 的某些子领域（比如只有几十道题的小子集），标准误差会大得多，排名就更加不稳定。

\begin{importantbox}{评估的核心哲学：没有``唯一正确''的评估}
这一讲最根本的信息不是某个 benchmark 好或不好，而是：\textbf{评估本质上是一个多目标、多视角、多利益方的近似问题}。不存在一个 benchmark 能对所有人、在所有场景下都给出``正确''答案。

好的评估实践不是追求一个完美的分数，而是：(1) 明确你的目标是什么；(2) 选择与目标最相关的度量；(3) 诚实地报告实验条件和局限；(4) 用多个互补的评估方式交叉验证。任何单一数字——无论是 MMLU 分数、ELO 排名还是 pass@k——都只是从一个角度看到的投影，不是全貌。
\end{importantbox}

